% Cover letter for submission to Transportation Research Part C: Emerging
% Technologies (Elsevier).
% NOTE before sending: verify whether the journal runs single- or double-blind
% review -- the manuscript currently carries author name/affiliation inline
% (main.tex), which matches single-blind convention; elsarticle.cls supports a
% `doubleblind` class option that suppresses author/affiliation from the
% rendered PDF if double-blind is required.
\documentclass[11pt,a4paper]{article}
\usepackage[margin=1in]{geometry}
\usepackage{parskip}
\usepackage{url}
\usepackage{microtype}
\pagestyle{empty}
\sloppy

\begin{document}

\noindent
Quang-Vinh Dang\\
British University Vietnam\\
Hung Yen, Vietnam\\
\url{vinh.dq4@buv.edu.vn}

\vspace{1em}
\noindent
To the Editors,\\
\emph{Transportation Research Part C: Emerging Technologies}

\vspace{1em}

Dear Editors,

I am submitting the manuscript ``Fair and Transferable Graph Reinforcement Learning for
Urban Rail Transit Network Expansion Across Asian Cities'' for consideration as a
regular research article in \emph{Transportation Research Part C}.

The paper addresses the Metro Network Expansion Problem (MNEP): given an existing rail
network and a limited construction budget, decide which stations or line segments to add
next. Three recent papers have each taken a piece of this idea and shown it works --
a GNN-based spatial state representation (MetroGNN), a fairness-aware multi-objective
reward decomposed via Tchebycheff scalarization, and a tabular reinforcement-learning
method that argues deep RL's overhead is unnecessary at this problem's scale -- but no
method combines the first two, and none has been tested from small, data-sparse emerging
metro systems to the largest networks in the world. We propose a GNN-based,
multi-objective actor-critic policy that combines a heterogeneous spatial-graph encoder
with a Tchebycheff-scalarized reward over demand, equity, and accessibility, condition
the policy itself on the sampled objective-weight vector so one trained model spans the
full efficiency-equity trade-off surface, and add a cross-city curriculum-transfer
mechanism -- pretraining on one city's network and warm-starting another's -- that
responds directly to the training-cost critique of GNN-based methods without abandoning
the GNN.

The paper's contribution is threefold: (1) the first MNEP method to combine GNN spatial
representation with multi-objective fairness decomposition in a single policy; (2) a
cross-city curriculum and sample-efficiency design (prioritized replay, reward shaping)
that directly answers the training-cost and carbon-footprint critique raised against
GNN-based MNEP methods; and (3) a rail-only scope discipline -- using each city's
existing bus network purely as an auxiliary input signal, never as an action -- that lets
the method extend, for the first time in this literature, to small, data-sparse emerging
metro systems (Hanoi, Ho Chi Minh City) that prior MNEP studies have not evaluated on.

The empirical results are reported as found rather than smoothed toward the hoped-for
conclusion. A five-seed sweep on Xi'an finds no reliable advantage for the full method
over either architectural ablation on demand, equity, or coverage, and the tabular RL
baseline still wins on raw demand under a matched training budget -- a negative result
for model complexity that we report plainly rather than obscure. Tracing an earlier,
more favorable-looking equity finding, we discovered GPU training non-determinism severe
enough to invert which architecture a metric appears to favor, traced it to two distinct
causes (non-deterministic scatter reduction in the GNN encoder and an unseeded
environment reset), fixed both, and verified deterministic re-runs -- a reproducibility
hazard we believe generalizes beyond our own implementation to other GNN-based methods
built on the same default scatter-reduction operations, and one we think is of
independent interest to this journal's readership regardless of this paper's own
architectural comparisons. On Ho Chi Minh City's structurally distinct, irregular
polygon graph, an encoder ablation collapses to near-zero demand where the full method
does not, and warm-starting from Xi'an's trained policy shows genuine positive transfer
at initialization and, under a reduced training budget, reaches non-trivial demand and
consistently higher coverage than training from scratch -- direct evidence for the
curriculum's sample-efficiency claim on the one city where we could test it. A separately
implemented radiation-model accessibility measure does not share the coverage proxy's
architectural sensitivity, a construct-validity finding we report rather than omit.

All code -- the GNN encoder, multi-objective policy, training and evaluation scripts for
both cities, and the full test suite -- is publicly available now, not embargoed until
publication, at \url{https://github.com/vinhqdang/metro_expansion}. Data sources are
public third-party releases (Michailidis et al.'s \texttt{mo-tndp}, the CityScope/MIT
\texttt{CSL\_HCMC} repository, and \texttt{GZUPA/subway-traffic-data-set}), cited and
documented in the manuscript's Data Availability statement.

The manuscript is original, has not been published previously, and is not under
consideration for publication elsewhere. I am the sole author and approve its
submission. I declare no competing financial or personal interests, and the work
received no specific grant from any funding agency.

Thank you for your consideration.

\vspace{1em}
\noindent
Yours sincerely,\\[1em]
Quang-Vinh Dang\\
British University Vietnam, Hung Yen, Vietnam

\end{document}
